\documentclass[11pt,a4paper]{article}
\usepackage[T1]{fontenc}
\usepackage[utf8]{inputenc}
\usepackage{newtxtext}
\usepackage{amsmath,amsthm,mathtools}
\usepackage{newtxmath}
\usepackage[margin=27mm,headheight=15pt]{geometry}
\usepackage{microtype,booktabs,array,tabularx}
\usepackage[dvipsnames]{xcolor}
\usepackage{enumitem,fancyhdr,etoolbox}
\usepackage{hyperref}
\usepackage{bookmark}
\hypersetup{colorlinks=true,linkcolor=MidnightBlue,citecolor=MidnightBlue,
 urlcolor=MidnightBlue,pdftitle={Global Fueter Primitives and Spherical Singularities},
 pdfauthor={Research manuscript prepared with ChatGPT},
 pdfsubject={Global inversion, affine periods, spherical principal parts, Mittag-Leffler realization}}
\urlstyle{same}
\numberwithin{equation}{section}
\newtheorem{theorem}{Theorem}[section]
\newtheorem{proposition}[theorem]{Proposition}
\newtheorem{lemma}[theorem]{Lemma}
\newtheorem{corollary}[theorem]{Corollary}
\theoremstyle{definition}
\newtheorem{definition}[theorem]{Definition}
\newtheorem{example}[theorem]{Example}
\theoremstyle{remark}
\newtheorem{remark}[theorem]{Remark}
\newcommand{\R}{\mathbb R}
\newcommand{\C}{\mathbb C}
\newcommand{\HH}{\mathbb H}
\newcommand{\HC}{\mathbb H_{\C}}
\newcommand{\SSS}{\mathbb S}
\newcommand{\Cp}{\C_+}
\newcommand{\cS}{\mathscr S}
\newcommand{\cM}{\mathscr M}
\newcommand{\cI}{\mathcal I}
\newcommand{\cD}{\mathcal D}
\newcommand{\lap}{\Delta_4}
\newcommand{\T}{\mathcal T}
\newcommand{\Aff}{\operatorname{Aff}_{\HH}}
\newcommand{\Res}{\operatorname{Res}}
\newcommand{\resaff}{\operatorname{res}^{\mathrm{aff}}}
\newcommand{\Rea}{\operatorname{Re}}
\newcommand{\Ima}{\operatorname{Im}}
\newcommand{\wind}{\operatorname{wind}}
\newcommand{\id}{\operatorname{id}}
\newcommand{\dd}{\,\mathrm d}
\newcommand{\norm}[1]{\lVert#1\rVert}
\newcommand{\abs}[1]{\lvert#1\rvert}
\newcommand{\qnorm}{\mathfrak n}
\newcommand{\supsph}[1]{\norm{#1}_{\mathrm{sph}}}
\setlist{itemsep=3pt,topsep=4pt}
\setlength{\parskip}{2.5pt}
\setlength{\emergencystretch}{2em}
\pagestyle{fancy}
\fancyhf{}
\fancyhead[L]{\small Global Fueter primitives and spherical singularities}
\fancyhead[R]{\small\thepage}
\renewcommand{\headrulewidth}{0.3pt}
\title{\vspace{-10mm}\textbf{Global Fueter Primitives\\and Spherical Singularities}\\[6pt]
\large Affine periods, complete principal parts, and\\a Mittag--Leffler realization theorem}
\author{Research manuscript prepared with ChatGPT}
\date{21 September 2026}
\begin{document}
\maketitle
\thispagestyle{empty}
\begin{abstract}
The local inverse of the quaternionic Fueter transform has an eight-real-dimensional
ambiguity, but a local primitive need not descend to the original domain. We make
this distinction explicit. For an axial monogenic function $g=P+IQ$, we construct
two closed quaternion-valued one-forms depending only on $P,Q$. Their periods are
necessary and sufficient for a single-valued slice-regular Fueter primitive, and
are exactly the two coefficients of its affine monodromy. Explicit logarithmic
generators realize every period vector. On a planar base with $m$ independent
holes this gives a split cokernel of real dimension $8m$, quantitative correction
bounds, and a continuous normalized inverse on the zero-period subspace. Reflected
generators extend smoothly through the real axis and give concrete counterexamples
to unrestricted same-domain surjectivity.

Around a missing nonreal quaternionic two-sphere, every axial monogenic function
has a unique normal form consisting of an affine-logarithmic residue pair, the
Fueter transform of a complexified-quaternionic Laurent principal part, and a
regular remainder. Singular classes with uniform growth $O(\rho^{-N})$ have real
dimension $8N$. We prove a sharp $o(\rho^{-1})$ removability theorem, compute the
exact leading growth constants, and identify the additional obstructions at
higher orders. Finally, arbitrary locally finite collections of finite spherical
principal parts and affine residues are realized by reducing the problem to the
classical Mittag--Leffler theorem for the second derivative of a multivalued
holomorphic primitive.
\end{abstract}
\noindent\textbf{Keywords.} Quaternionic analysis; inverse Fueter mapping;
axially monogenic functions; period obstruction; affine monodromy;
spherical singularity; Mittag--Leffler theorem.

\medskip
\noindent\textbf{Status and scope.} This is a research continuation of the supplied
exposition~\cite{Source}, not a claim to have discovered the local Fueter inverse.
The proposed contributions are the explicit global obstruction, splitting,
normal-form, sharp-growth, and realization results proved below. We did not locate
this package of statements in the primary sources checked for this manuscript;
that is not a guarantee of priority. The proofs have been supplemented by symbolic
and numerical checks, but have not been independently refereed or formally verified.
The problem settled here is a precise global-inversion and singularity problem,
not a claim to settle every open question in quaternionic analysis.
\newpage
\begingroup
\small
\makeatletter
\patchcmd{\l@section}{1.0em}{0.55em}{}{}
\makeatother
\tableofcontents
\endgroup
\newpage

\section{The problem and the main results}
\label{sec:intro}

The supplied article~\cite{Source} develops slice regularity and Cauchy--Fueter
regularity as complementary theories and ends with a constructive axial inverse
of the Laplacian on a simply connected upper-half-plane base. That construction
integrates two closed one-forms in succession. It also explicitly notes that
periods may prevent globalization. We take that qualification as our starting
point, keeping its conventions: left regularity, quaternionic coefficients on
the right, and a central auxiliary complex unit $\iota$.

Local inverse Fueter theory is established mathematics
\cite{CSS,CPSS}. It is therefore not enough to integrate once more on a disk and
call the result new. The substantive questions are these: can global
solvability be tested directly from the target function, without first choosing
local primitives? Which period vectors actually occur? How do these global
defects interact with singularities on an entire conjugacy sphere? Can arbitrary
singular data be realized simultaneously without leaving the axial monogenic
class?

Here are the answers, with the analytic setting defined precisely in
Section~\ref{sec:setting}. If $U\subset\Cp$ is connected and
$g(x+Ir)=P(x,r)+IQ(x,r)$ is axial monogenic, set
\begin{equation}
 \omega_g=\frac{-P\dd x+Q\dd r}{2},\qquad
 \theta_g=\frac{(xP+rQ)\dd x+(rP-xQ)\dd r}{2}.
 \label{eq:forms-intro}
\end{equation}
Both forms are closed. A single-valued slice-regular $f$ with $\lap f=g$ exists
on the original circularized domain if and only if \emph{both} forms have zero
periods. For a loop $\gamma$, their periods $(a_\gamma,b_\gamma)$ are precisely
the coefficients in the change
\[
 f(q)\longmapsto f(q)+q a_\gamma+b_\gamma
\]
under analytic continuation. This criterion is proved in
Theorem~\ref{thm:period}; it does not assume finite connectivity.

When $U$ has a finite winding-number basis of $m$ holes, explicit generators
produce the exact sequence
\begin{equation}
 0\longrightarrow\HH^2\longrightarrow\cS(U)
 \xrightarrow{\ \T\ }\cM(U)
 \xrightarrow{\ \Pi\ }(\HH^2)^m\longrightarrow0.
 \label{eq:sequence-intro}
\end{equation}
Here $\cS(U)$ is the space of holomorphic quaternionic stems on $U$,
$\cM(U)$ is the space of axial monogenic functions, and $\T F=\lap\cI F$.
The sequence is one of right $\HH$-modules, or equivalently of real vector
spaces. In particular, the obstruction has \emph{real dimension $8m$}, not
merely $m$ or $4m$. Sections~\ref{sec:generators}--\ref{sec:stability} provide
an explicit complement, norm estimates, and normalized inversion.

The singularity results are not just a reformulation of this dimension count.
At a nonreal sphere $S_p$, let $\rho$ denote distance to that sphere.
Theorems~\ref{thm:normal} and~\ref{thm:growth} give a complete singular normal
form and show that the space of $O(\rho^{-N})$ singular germs modulo removable
germs has real dimension $8N$. Theorem~\ref{thm:asymptotic} calculates the exact
leading constants. The critical exponent is $1$, and at that exponent an
eight-real-dimensional affine residue pair detects removability completely.
At higher exponents it does not.

Finally, Theorem~\ref{thm:ML} realizes arbitrary finite principal parts and
residue pairs at any locally finite family of spheres over a simply connected
upper base. The proof uses the fact that the second derivative of a local
holomorphic Fueter primitive is globally single-valued even when the primitive
itself is not.

\subsection{A necessary distinction about surjectivity}
\label{subsec:surjectivity-context}

Statements called inverse or surjectivity theorems occur with different domain
and continuation conventions. The explicit inverse theorem of
Colombo--Pe\~na Pe\~na--Sabadini--Sommen~\cite[Theorem~3]{CPSS} is formulated on
a rectangular axial base. Dong--Qian~\cite[Theorem~3.6]{DQ} give a
contour-based surjectivity formulation. Neither the word ``inverse'' nor a
local integral formula by itself removes monodromy.

There is also an explicitly unqualified same-domain surjectivity sentence in
Krau\ss har--Perotti~\cite[\S4, p.~2539, before Proposition~24]{KP}, after the
definition of axially monogenic functions on an axially symmetric domain.
The examples in Section~\ref{sec:realaxis} show that this unrestricted
formulation needs a topological qualification when functions are required to
be single-valued on the unchanged domain. This observation does not purport
to refute that paper's polynomial calculations or its particular eigenvalue
conclusions. Our replacement is a necessary-and-sufficient period criterion,
proved directly, together with explicit obstructed functions.

\section{Analytic setting and the local identities}
\label{sec:setting}

\subsection{Upper stems and circularization}
Write
\[
 \HH=\R+\mathbf i\R+\mathbf j\R+\mathbf k\R,
 \qquad \SSS=\{I\in\HH:I^2=-1\},\qquad
 \Cp=\{x+\iota r:r>0\}.
\]
The auxiliary unit $\iota$ is central in
$\HC=\HH\oplus\iota\HH$. Complex conjugation on this algebra is
$\kappa(A+\iota B)=A-\iota B$; it is not quaternionic conjugation.
For a connected open $U\subset\Cp$, put
\[
 \Omega_U=\{x+Ir:x+\iota r\in U,\ I\in\SSS\}.
\]
The map $(z,I)\mapsto x+Ir$ identifies this domain smoothly with
$U\times\SSS$.

Let
\[
 \cS(U)=\mathcal O(U,\HC).
\]
For $F=A+\iota B\in\cS(U)$, where $A,B:U\to\HH$, define
\begin{equation}
 \cI F(x+Ir)=A(x,r)+IB(x,r).
 \label{eq:slice}
\end{equation}
This is the usual stem construction of a left slice-regular function on a
product domain: extend $F$ to $\bar U$ by
$F(\bar z)=\kappa(F(z))$. Its Cauchy--Riemann equations are
\begin{equation}
 A_x=B_r,\qquad A_r=-B_x.
 \label{eq:CR}
\end{equation}
The general stem formalism is due to Ghiloni--Perotti~\cite{GP}.
All holomorphic and meromorphic statements involving $\HC$ below mean the
corresponding statements in four complex coordinates.

The letters $a,b$ denote quaternions unless explicitly stated otherwise.
A coefficient $c\in\HC$ may have the form $c=\alpha+\iota\beta$ with
$\alpha,\beta\in\HH$. Keeping these two coefficient spaces separate is
essential. In particular, $\T$ below is right $\HH$-linear but is \emph{not}
complex-linear with respect to $\iota$.

\subsection{The axial monogenic space}
The left Cauchy--Fueter operator and Euclidean Laplacian are
\[
 \cD=\partial_x+\mathbf i\partial_{x_1}
     +\mathbf j\partial_{x_2}+\mathbf k\partial_{x_3},
 \qquad \lap=\partial_x^2+\partial_{x_1}^2
                  +\partial_{x_2}^2+\partial_{x_3}^2.
\]
An \emph{axial monogenic function} is a smooth function
\begin{equation}
 g(x+Ir)=P(x,r)+IQ(x,r),\qquad P,Q:U\to\HH,
 \label{eq:axial}
\end{equation}
with $\cD g=0$. Denote its space by $\cM(U)$. The values at $I$ and $-I$
recover $P,Q$ uniquely, so this definition is independent of a preferred
quaternionic slice.

\begin{lemma}[Axial identities and the Fueter transform]
\label{lem:identities}
For a smooth axial function $A+IB$,
\begin{align}
 \cD(A+IB)&=A_x-B_r-\frac{2B}{r}+I(B_x+A_r),
 \label{eq:axialD}\\
 \lap(A+IB)&=A_{xx}+A_{rr}+\frac{2A_r}{r}
  +I\left(B_{xx}+B_{rr}+\frac{2B_r}{r}-\frac{2B}{r^2}\right).
 \label{eq:axialDelta}
\end{align}
Consequently $g=P+IQ$ belongs to $\cM(U)$ precisely when
\begin{equation}
 P_x-Q_r-\frac{2Q}{r}=0,\qquad Q_x+P_r=0.
 \label{eq:Vekua}
\end{equation}
For $F=A+\iota B\in\cS(U)$, the transform $\T F=\lap\cI F$ is given by
\begin{equation}
 \T F=P+IQ,\qquad
 P=\frac{2A_r}{r},\qquad
 Q=\frac{2B_r}{r}-\frac{2B}{r^2},
 \label{eq:T}
\end{equation}
and belongs to $\cM(U)$.
\end{lemma}
\begin{proof}
Put $v=\mathbf i x_1+\mathbf jx_2+\mathbf kx_3$ and $I=v/r$.
For $\ell=1,2,3$ and $e_1=\mathbf i,e_2=\mathbf j,e_3=\mathbf k$,
\[
 \partial_{x_\ell}r=\frac{x_\ell}{r},\qquad
 \partial_{x_\ell}I=\frac{e_\ell}{r}-\frac{v x_\ell}{r^3}.
\]
Thus $\sum e_\ell\partial_{x_\ell}I=-2/r$,
$\Delta_{\R^3}I=-2I/r^2$, and
$\sum(\partial_{x_\ell}I)(\partial_{x_\ell}r)=0$.
The product rule gives \eqref{eq:axialD}--\eqref{eq:axialDelta}.
An expression $X+IY$ that vanishes for every $I\in\SSS$ has $X=Y=0$;
this proves \eqref{eq:Vekua}.

For a holomorphic stem, \eqref{eq:CR} implies that $A,B$ are planar
harmonic, giving \eqref{eq:T}. Direct differentiation then yields
\[
 P_x=\frac{2B_{rr}}r=Q_r+\frac{2Q}r,
 \qquad
 Q_x+P_r=\frac{2(B_{rx}+A_{rr})}r
             -\frac{2(B_x+A_r)}{r^2}=0.
\]
This proves the mapping assertion.
\end{proof}

These are established local identities, also used in the supplied
exposition~\cite{Source}; the affine kernel is recorded in
\cite[Lemma~23]{KP}. We reproduce their proofs because the global arguments
need their exact signs and coefficient conventions.

\begin{lemma}[Two coordinates of a primitive]
\label{lem:coordinates}
If $F=A+\iota B\in\cS(U)$ and $g=\T F$, define quaternion-valued functions
\begin{equation}
 C=\frac Br,\qquad E=A-x\frac Br.
 \label{eq:CE}
\end{equation}
Then
\begin{equation}
 dC=\omega_g,\qquad dE=\theta_g,
 \label{eq:CEforms}
\end{equation}
where the forms are those in \eqref{eq:forms-intro}. Moreover,
\begin{equation}
 F=zC+E,\qquad \cI F(q)=qC+E.
 \label{eq:reconstruct}
\end{equation}
\end{lemma}
\begin{proof}
Equations \eqref{eq:CR} and \eqref{eq:T} give
\[
 C_x=-\frac P2,\quad C_r=\frac Q2,\quad
 E_x=B_r-C+\frac{xP}{2}=\frac{xP+rQ}{2},\quad
 E_r=\frac{rP-xQ}{2}.
\]
These are exactly \eqref{eq:CEforms}. Substitution proves
\eqref{eq:reconstruct}.
\end{proof}

\begin{corollary}[Kernel, including product domains]
\label{cor:kernel}
On every connected $U\subset\Cp$,
\[
 \ker\T=\Aff(U):=\{z\mapsto za+b:a,b\in\HH\}.
\]
\end{corollary}
\begin{proof}
If $\T F=0$, Lemma~\ref{lem:coordinates} makes $C,E$ constant on $U$;
then $F=za+b$. The converse follows from \eqref{eq:T}.
\end{proof}

\begin{remark}[Why complex-affine is not the kernel]
\label{rem:complexaffine}
The holomorphic stem $F=\iota$ induces $f(x+Ir)=I$, and
$\T F=-2I/r^2\ne0$. More generally,
\begin{equation}
 \T\bigl(\iota(za+b)\bigr)
 =-\frac{2a}{r}-\frac{2I(xa+b)}{r^2}.
 \label{eq:imag-affine}
\end{equation}
Thus replacing $\HH$ by $\HC$ in the kernel statement would erase the
very obstruction that the subsequent theory measures.
\end{remark}

\section{The complete global period obstruction}
\label{sec:periods}

\begin{lemma}[Closedness without a primitive]
\label{lem:closed}
For every $g\in\cM(U)$, both $\omega_g$ and $\theta_g$ are closed.
\end{lemma}
\begin{proof}
Using the orientation $dx\wedge dr$,
\[
 2d\omega_g=(Q_x+P_r)\,dx\wedge dr=0,
\]
while
\[
 2d\theta_g=
 \bigl(r(P_x-Q_r)-2Q-x(Q_x+P_r)\bigr)\,dx\wedge dr=0
\]
by \eqref{eq:Vekua}. The calculation is componentwise over $\R$.
\end{proof}

\begin{theorem}[Global primitive criterion and affine monodromy]
\label{thm:period}
Let $U\subset\Cp$ be any connected open set and let $g\in\cM(U)$.
For a piecewise $C^1$ closed curve $\gamma\subset U$, define
\begin{equation}
 a_\gamma(g)=\int_\gamma\omega_g,\qquad
 b_\gamma(g)=\int_\gamma\theta_g.
 \label{eq:periods}
\end{equation}
The following conditions are equivalent:
\begin{enumerate}[label=\textup{(\roman*)}]
\item There is a single-valued $F\in\cS(U)$ with $\T F=g$.
\item Both one-forms are exact on $U$.
\item $a_\gamma(g)=b_\gamma(g)=0$ for every closed curve $\gamma$.
\end{enumerate}
If these conditions hold, all primitives are
\begin{equation}
 F(z)=z\left(a+\int_{z_*}^{z}\omega_g\right)
       +b+\int_{z_*}^{z}\theta_g,
 \qquad a,b\in\HH,
 \label{eq:globalinverse}
\end{equation}
with path-independent integrals.

Without any vanishing assumption, a primitive exists on the universal cover.
Analytic continuation along a loop $\gamma$ changes that primitive by
\begin{equation}
 F(z)\longmapsto F(z)+z a_\gamma(g)+b_\gamma(g).
 \label{eq:monodromy}
\end{equation}
The period pair defines an additive homomorphism
$H_1(U;\mathbb Z)\to\HH^2$.
\end{theorem}
\begin{proof}
Necessity of exactness is Lemma~\ref{lem:coordinates}. Exact forms have zero
periods. Conversely, when periods vanish, integrating each of the four real
components of the closed forms from $z_*$ defines globally single-valued
$C,E$ with $dC=\omega_g$ and $dE=\theta_g$. Path independence follows by
subtracting the two paths to obtain a closed curve. Set
\[
 A=xC+E,\qquad B=rC.
\]
The derivatives of the primitives give
\[
 A_x=C+\frac{rQ}{2}=B_r,\qquad
 A_r=\frac{rP}{2}=-B_x.
\]
Therefore $F=A+\iota B=zC+E$ is holomorphic. Formula \eqref{eq:T} gives
$\T F=g$. The constants of integration produce exactly
\eqref{eq:globalinverse}; Corollary~\ref{cor:kernel} proves that there is
no additional ambiguity.

On the universal cover the pulled-back closed forms have primitives. A deck
continuation adds their periods to $C,E$, and hence adds
$za_\gamma+b_\gamma$ to $F$. These changes are themselves in the kernel
of $\T$. The integral of a closed form is unchanged by a homology, by
Stokes' theorem applied to each component; concatenation adds integrals.
Thus the pair factors through $H_1(U;\mathbb Z)$. Finally, a nonzero pair
cannot give the identically zero affine stem: the $\iota$-component of
$za+b$ is $ra$, so $za+b=0$ on an open upper set forces $a=b=0$.
\end{proof}

\subsection{Relation to the two successive integrations}
The local construction in~\cite{Source} first obtains $C$ from $\omega_g$,
then integrates
\[
 \eta=\left(C+\frac{rQ}{2}\right)dx+\frac{rP}{2}dr
\]
to obtain $A$. Since
\begin{equation}
 \eta-d(xC)=\theta_g,
 \label{eq:eliminateC}
\end{equation}
the second obstruction is already present in $g$ itself. It does not depend
on the chosen value of $C$ at the base point, nor must it be recomputed
for different local primitives. This elimination is the key that turns the
local construction into a simultaneous global test.

\begin{remark}[Additive, not noncommutative, monodromy]
Although the target algebra is noncommutative, the monodromy here is the
additive group of affine functions with right quaternionic coefficients.
No path-ordered exponential or multiplicative quaternionic holonomy occurs.
The distinction matters when counting the obstruction or computing it
numerically.
\end{remark}

\begin{proposition}[Changing the real origin]
\label{prop:translation}
Under $x'=x-c$, $r'=r$ for a real constant $c$, the same target function has
periods
\[
 a'_\gamma=a_\gamma,\qquad b'_\gamma=b_\gamma+c a_\gamma.
\]
For a sphere centered at the real coordinate $u$, the pair
$(a,ua+b)$ is therefore unchanged by a simultaneous translation of coordinates.
\end{proposition}
\begin{proof}
The first form is unchanged, while direct substitution gives
$\theta'_g=\theta_g+c\omega_g$. The sphere center becomes $u-c$.
\end{proof}

\section{Explicit generators and a split global obstruction space}
\label{sec:generators}

\subsection{Logarithmic primitives with single-valued transforms}
Let $p=u+\iota v\in\C\setminus U$; $v$ need not be positive for the formula
below. On any local branch put
\begin{equation}
 F_{p,a,b}(z)=\frac{za+b}{2\pi\iota}\log(z-p),
 \qquad a,b\in\HH.
 \label{eq:logprimitive}
\end{equation}
Changing the branch of the logarithm by $2\pi\iota n$ adds
$n(za+b)$. Its Fueter transform is thus single-valued on $U$; denote it by
$G_{p,a,b}$.

\begin{proposition}[Branch-free formula and normalized periods]
\label{prop:generators}
Let $R_p=(x-u)^2+(r-v)^2$ and $s_p=\tfrac12\log R_p$.
Then $G_{p,a,b}=P_{p,a,b}+IQ_{p,a,b}$, where
\begin{align}
 P_{p,a,b}&=\frac1\pi\left(
 \frac{(r-v)a}{R_p}+\frac{s_p a}{r}
 +\frac{(x-u)(xa+b)}{rR_p}\right),
 \label{eq:G-P}\\
 Q_{p,a,b}&=\frac1\pi\left(
 \frac{(x-u)a}{R_p}+\frac{s_p(xa+b)}{r^2}
 -\frac{(r-v)(xa+b)}{rR_p}\right).
 \label{eq:G-Q}
\end{align}
For every closed curve $\gamma\subset U$,
\begin{equation}
 \left(\int_\gamma\omega_{G_{p,a,b}},
       \int_\gamma\theta_{G_{p,a,b}}\right)
 =\wind(\gamma,p)(a,b).
 \label{eq:G-period}
\end{equation}
\end{proposition}
\begin{proof}
Write $\log(z-p)=s_p+\iota t$ on a local branch. If $L=xa+b$,
\[
 A=\frac{Lt+ra s_p}{2\pi},\qquad
 B=\frac{ra t-Ls_p}{2\pi},\qquad
 C=\frac{a t-Ls_p/r}{2\pi}.
\]
Since
\[
 t_x=-\frac{r-v}{R_p},\quad t_r=\frac{x-u}{R_p},\quad
 (s_p)_x=\frac{x-u}{R_p},\quad (s_p)_r=\frac{r-v}{R_p},
\]
the identities $P=-2C_x$, $Q=2C_r$ give
\eqref{eq:G-P}--\eqref{eq:G-Q}. These expressions use only the real logarithm
of a positive number; there is no hidden choice of argument. Local Fueter
mapping proves that they are monogenic everywhere on $U$.

Continuing \eqref{eq:logprimitive} along $\gamma$ changes it by
$\wind(\gamma,p)(za+b)$. Comparing with the uniquely determined monodromy
in Theorem~\ref{thm:period} proves \eqref{eq:G-period}.
\end{proof}

\subsection{Finite connectivity, with explicit topological hypotheses}
\begin{definition}[A finite winding basis]
\label{def:finitebasis}
We say that $U$ has a finite winding basis of size $m$ if there exist points
$p_1,\ldots,p_m\in\Cp\setminus U$ and piecewise $C^1$ cycles
$\gamma_1,\ldots,\gamma_m\subset U$ whose classes form a $\mathbb Z$-basis
of $H_1(U;\mathbb Z)$ and satisfy
\[
 \wind(\gamma_j,p_k)=\delta_{jk}.
\]
The case $m=0$ means $H_1(U;\mathbb Z)=0$.
\end{definition}
This hypothesis is satisfied, for example, by an upper disk with finitely many
disjoint closed disks removed from its interior, or by a simply connected
upper domain punctured at finitely many points. It states explicitly the
topological property used in the proof; no boundary smoothness beyond the
chosen cycles is needed.

\begin{theorem}[Realized cokernel and canonical correction after choices]
\label{thm:splitting}
Suppose $U$ has a finite winding basis of size $m$. Define
\[
 \Pi g=\bigl(a_{\gamma_j}(g),b_{\gamma_j}(g)\bigr)_{j=1}^m,
 \qquad
 S\bigl((a_j,b_j)_j\bigr)=\sum_{j=1}^mG_{p_j,a_j,b_j}.
\]
Then \eqref{eq:sequence-intro} is exact, $\Pi S=\id$, and
\begin{equation}
 \cM(U)=\T\cS(U)\oplus S((\HH^2)^m).
 \label{eq:directsum}
\end{equation}
The operator $R=\id-S\Pi$ is a projection onto the space of targets having
a single-valued global Fueter primitive. In particular,
\[
 \dim_\R\bigl(\cM(U)/\T\cS(U)\bigr)=8m.
\]
The complement depends on the selected points and cycles, but the quotient
and the criterion for belonging to $\T\cS(U)$ do not.
\end{theorem}
\begin{proof}
The injection of $\HH^2$ is $(a,b)\mapsto za+b$. Its image is the
kernel by Corollary~\ref{cor:kernel}. By the winding normalization and
\eqref{eq:G-period}, $\Pi S=\id$, proving surjectivity of $\Pi$.
Theorem~\ref{thm:period}, together with the homology-basis hypothesis, gives
$\ker\Pi=\T\cS(U)$. For any $g$,
\[
 g=(g-S\Pi g)+S\Pi g,
\]
and the first summand is in $\ker\Pi$. If $S\mathbf v$ also lies in that
kernel, then $\mathbf v=\Pi S\mathbf v=0$. This proves the direct sum.
Finally $\HH^2$ has real dimension eight.
\end{proof}

\begin{example}[The obstructions are independent]
\label{ex:independent}
On a punctured upper disk containing $p$ in the missing point,
$G_{p,1,0}$ has first period $1$ and second period $0$.
The function $G_{p,0,1}$ has first period $0$ and second period $1$.
Neither has a single-valued Fueter primitive. Thus checking only the form
used in the first integration is insufficient, and so is checking only the
second form. Right multiplication by a quaternion produces four real
independent directions in each of these two types.
\end{example}

\section{Quantitative correction, closed range, and stable inversion}
\label{sec:stability}

For a compact planar set $K\subset U$, use the actual quaternionic supremum
\begin{equation}
 \norm{g}_K=\sup_{z=x+\iota r\in K}\ \sup_{I\in\SSS}|P(x,r)+IQ(x,r)|.
 \label{eq:normK}
\end{equation}
In particular, norms involve whole conjugacy spheres, not just one section.
The identity
\[
 \frac{|P+IQ|^2+|P-IQ|^2}{2}=|P|^2+|Q|^2
\]
shows that $(|P|^2+|Q|^2)^{1/2}\le\sup_I|P+IQ|$.

\begin{proposition}[Period bounds and distance from the exact subspace]
\label{prop:bounds}
If a cycle $\gamma$ has length $\ell$ and
$R_\gamma=\max_{z\in\gamma}|z|$, then
\begin{equation}
 |a_\gamma(g)|\le\frac{\ell}{2}\norm{g}_\gamma,
 \qquad
 |b_\gamma(g)|\le\frac{R_\gamma\ell}{2}\norm{g}_\gamma.
 \label{eq:periodbounds}
\end{equation}
For every $h\in\T\cS(U)$,
\begin{equation}
 \norm{g-h}_\gamma\ge
 \max\left\{\frac{2|a_\gamma(g)|}{\ell},
             \frac{2|b_\gamma(g)|}{R_\gamma\ell}\right\}.
 \label{eq:distancebound}
\end{equation}
\end{proposition}
\begin{proof}
For a real tangent $(\dot x,\dot r)$, Cauchy--Schwarz in a real Euclidean
space gives
\[
 |\omega_g(\dot x,\dot r)|
 \le\frac12(|P|^2+|Q|^2)^{1/2}
                    (\dot x^2+\dot r^2)^{1/2}.
\]
The coefficient vector of $P,Q$ in $2\theta_g(\dot x,\dot r)$ is
$(x\dot x+r\dot r,r\dot x-x\dot r)$, whose squared norm is
$|z|^2(\dot x^2+\dot r^2)$. Integrating proves
\eqref{eq:periodbounds}. Apply those inequalities to $g-h$ and use the
vanishing periods of $h$ to get \eqref{eq:distancebound}.
\end{proof}

\begin{corollary}[Closed range and bounded finite-dimensional correction]
\label{cor:closedrange}
In the topology of locally uniform convergence on circularized compact sets,
$\T\cS(U)$ is closed relative to $\cM(U)$, even when $U$ has infinitely
many holes. Under the finite-basis hypothesis, $S$, $\Pi$, and $R$ are
continuous. More explicitly,
\begin{equation}
 \norm{S\Pi g}_K\le\sum_{j=1}^m
 \left(|a_{\gamma_j}(g)|\,\norm{G_{p_j,1,0}}_K
      +|b_{\gamma_j}(g)|\,\norm{G_{p_j,0,1}}_K\right).
 \label{eq:correctionbound}
\end{equation}
\end{corollary}
\begin{proof}
Each period is a continuous functional by \eqref{eq:periodbounds}.
The exact subspace is the intersection of all their kernels, by
Theorem~\ref{thm:period}, and is therefore closed. Formula
\eqref{eq:correctionbound} follows from right $\HH$-linearity, the triangle
inequality, and $|uv|=|u||v|$ for quaternions. It implies the remaining
continuity assertions with the finite collection of period estimates.
\end{proof}

\begin{theorem}[Normalized inverse and a pathwise estimate]
\label{thm:stableinverse}
Fix $z_*\in U$. For every zero-period $g\in\cM(U)$ there is a unique
primitive $F\in\cS(U)$ satisfying $F(z_*)=0$.
Suppose a path from $z_*$ to $z$ has length at most $L$, its points have
modulus at most $R_0$, and $|g|\le M$ on all spheres over that path.
Then
\begin{equation}
 |C(z)|\le\frac{LM}{2},\qquad
 |E(z)|\le\frac{R_0LM}{2},\qquad
 |\cI F(q)|\le\frac{(|q|+R_0)LM}{2}
 \quad(q=x+Ir).
 \label{eq:inversebound}
\end{equation}
The normalized inverse is continuous in the corresponding compact-open
topologies.
\end{theorem}
\begin{proof}
Take $C(z_*)=E(z_*)=0$ in \eqref{eq:globalinverse}. The pointwise
one-form estimates used in Proposition~\ref{prop:bounds} give the first
two inequalities; $\cI F=qC+E$ gives the third. If an affine kernel
function $za+b$ vanishes at $z_*=x_*+\iota r_*$, its imaginary component
$r_*a$ vanishes. Since $r_*>0$, $a=b=0$, proving uniqueness.

For completeness, any compact $K\Subset U$ admits a uniform version of
the path hypothesis on a slightly larger compact set. Cover $K$ by
finitely many disks with closures in $U$, join their centers to $z_*$ by
finitely many polygonal paths in $U$, and then join points of each disk to
its center. The union is contained in a compact $K'\Subset U$, and the
path lengths are uniformly bounded. Apply \eqref{eq:inversebound} on
$K'$ to differences of targets. The complexified norm
$(|A|^2+|B|^2)^{1/2}$ of $F$ is bounded in the same way by
$|z||C|+|E|$. Thus the stems converge locally uniformly as well.
\end{proof}

The finite correction is not an unspecified existence device. The procedure
is to integrate the two explicit forms on a homology basis, subtract the
corresponding sum of \eqref{eq:G-P}--\eqref{eq:G-Q}, and integrate the
remaining exact forms along any convenient paths. Formula
\eqref{eq:distancebound} also shows that a genuine period defect cannot be
made arbitrarily small on its detecting loop by choosing better exact targets.

\section{Crossing the real axis and an explicit failure of global surjectivity}
\label{sec:realaxis}

The preceding domains avoid the real axis, but the obstruction is not an
artifact of that restriction. We now treat genuine axially symmetric slice
domains.

Let $D\subset\C$ be connected, open, conjugation invariant, and intersecting
$\R$. Assume that $U=D\cap\Cp$ is connected, and let $\Omega_D$ be its
usual circularization, including the real points. A holomorphic stem on $D$
satisfies $F(\bar z)=\kappa(F(z))$. Suppose $g$ is a smooth axial monogenic
function on $\Omega_D$.

\begin{theorem}[The same period criterion across the real axis]
\label{thm:crossreal}
There is a global slice-regular $f$ on $\Omega_D$ with $\lap f=g$ if and
only if the two forms on $U$ have zero periods. Such a primitive is unique
up to $qa+b$, $a,b\in\HH$.
\end{theorem}
\begin{proof}
Necessity follows by restriction to $\Omega_U$. For sufficiency, construct
$C,E$ on $U$ by Theorem~\ref{thm:period}. We must prove that
$A=xC+E$ and $B=rC$ have the required extension at every real interval in $D$.

Fix an imaginary unit $I_0$. From the values of $g$ on its two half-slices,
\[
 P(x,r)=\frac{g(x+I_0r)+g(x-I_0r)}2,\qquad
 Q(x,r)=-\frac{I_0}{2}\bigl(g(x+I_0r)-g(x-I_0r)\bigr).
\]
These extend smoothly to signed $r$, with $P$ even and $Q$ odd. The two
one-forms, which have no divisions by $r$, therefore extend smoothly across
the real axis. Their exterior derivatives still vanish: the Vekua identities
hold for $r\ne0$, and their denominator-free consequences in
Lemma~\ref{lem:closed} extend by continuity.

On a small disk centered at a real point, choose primitives of these extended
forms. Adjust their constants to agree with $C,E$ in the upper half-disk.
Reflection $r\mapsto-r$ preserves each form. The reflected primitive and
the original differ by a constant, which is zero at a real point; hence both
local primitives are even in $r$. It follows that $A$ is even, $B$ is odd,
and their Cauchy--Riemann equations hold through the real line. Thus
$F=A+\iota B$ is holomorphic there. Extending the global upper stem to the
lower half by reflection now gives a holomorphic stem on all of $D$.
The local constructions agree because they agree on their upper portions.
The equality $\lap f=g$ extends across the real line by continuity.
Uniqueness follows from the affine kernel on the connected upper base.
\end{proof}

\begin{proposition}[Reflected generators]
\label{prop:reflected}
For $p\in\Cp$ and $a,b\in\HH$, set on $r>0$
\begin{equation}
 G^{\mathrm{sym}}_{p,a,b}=G_{p,a,b}-G_{\bar p,a,b}.
 \label{eq:reflectedG}
\end{equation}
This function extends to a smooth axial monogenic function on
$\HH\setminus S_p$, where
\[
 S_p=u+v\SSS,\qquad p=u+\iota v.
\]
Its periods on an upper loop are
$\wind(\gamma,p)(a,b)$. Consequently, if $U=D\cap\Cp$ has a finite
winding basis of size $m$, the quotient of smooth axial monogenic functions
on $\Omega_D$ by global Fueter transforms has real dimension $8m$.
\end{proposition}
\begin{proof}
A local primitive is
\begin{equation}
 F^{\mathrm{sym}}_{p,a,b}(z)
  =\frac{za+b}{2\pi\iota}
       \log\frac{z-p}{z-\bar p}.
 \label{eq:reflectedF}
\end{equation}
Branch changes add $n(za+b)$, so its transform is single-valued. The ratio
inside the logarithm is nonzero and holomorphic away from $p,\bar p$.
Near a real point it has modulus one on the real line. A local logarithm
can therefore be chosen purely imaginary on a real interval. With that
choice the function in \eqref{eq:reflectedF} is a reflection-compatible
holomorphic stem. Its induced slice function is real analytic through the
real axis; so is its Fueter transform. These local transforms agree with
\eqref{eq:reflectedG} and with one another, since their primitive differences
are real-affine.

A loop contained in $\Cp$ has winding number zero about $\bar p$.
Proposition~\ref{prop:generators} gives the asserted periods. The last
statement follows from Theorem~\ref{thm:crossreal} and the same splitting
argument as in Theorem~\ref{thm:splitting}, using the reflected generators.
\end{proof}

\subsection{Two normalized examples on the complement of the unit sphere}
\label{subsec:counterexample}
Take
\[
 D=\C\setminus\{\iota,-\iota\},\qquad
 \Omega_D=\HH\setminus\SSS.
\]
Consider the local scalar holomorphic function
\begin{equation}
 h(z)=\frac{1}{2\pi\iota}\log\frac{z-\iota}{z+\iota}.
 \label{eq:hcounter}
\end{equation}
Its derivative is the globally single-valued rational function
\begin{equation}
 h'(z)=\frac1{\pi(1+z^2)}.
 \label{eq:hprime}
\end{equation}
The increment of $h$ around a positively oriented loop about $\iota$ is
$1$. This is also seen by integrating \eqref{eq:hprime}, whose residue at
$\iota$ is $1/(2\pi\iota)$.

\begin{corollary}[Two concrete same-domain obstructions]
\label{cor:counterexample}
The functions
\[
 g_b=\T h=G^{\mathrm{sym}}_{\iota,0,1},\qquad
 g_a=\T(zh)=G^{\mathrm{sym}}_{\iota,1,0}
\]
are globally smooth axial monogenic functions on $\HH\setminus\SSS$,
but neither is the Laplacian of any globally single-valued slice-regular
function on that same domain. Their upper-loop period pairs are respectively
$(0,1)$ and $(1,0)$.
\end{corollary}
\begin{proof}
Smoothness and monogenicity follow from Proposition~\ref{prop:reflected}.
The monodromy of the local primitives is respectively $1$ and $z$.
Theorem~\ref{thm:period} identifies the corresponding period pairs and
rules out a single-valued primitive. Equivalently, if a global primitive
existed, its difference from the displayed local one would be $za+b$ with
fixed quaternionic coefficients on the universal cover. Such a difference
has no monodromy and cannot cancel $1$ or $z$.
\end{proof}

These are branch-free target functions, not multivalued functions disguised
by notation. For example, $g_b=-2(C_b)_x+2I(C_b)_r$ with the globally
single-valued real function
\begin{equation}
 C_b(x,r)=-\frac1{2\pi r}
       \log\frac{|x+\iota r-\iota|}{|x+\iota r+\iota|}.
 \label{eq:Cb}
\end{equation}
The apparent singularity at $r=0$ is removable, with
$C_b(x,0)=1/(\pi(1+x^2))$. Thus even the first potential can be globally
single-valued while the second potential cannot. The real-axis values are
particularly simple:
\begin{equation}
 g_b(x)=\frac{4x}{\pi(1+x^2)^2},\qquad
 g_a(x)=-\frac4{\pi(1+x^2)^2}.
 \label{eq:realtraces}
\end{equation}
To verify them, note that for a reflection-compatible scalar holomorphic
stem $H$, the real-axis value of $\T H$ is $-2H''(x)$. This follows from
\eqref{eq:T} and its Taylor expansion at $r=0$. Apply it to $h$ and $zh$,
using \eqref{eq:hprime}.

There is a second, purely complex-analytic check of the obstruction for
$g_b$. Suppose that a global reflection-compatible holomorphic stem $F$ on
$D$ satisfied $\T F=g_b$. The real-axis identity just used gives
\[
 F''(x)=-\frac{2x}{\pi(1+x^2)^2}.
\]
The componentwise holomorphic identity theorem extends this equality to
all of $D$. Consequently
\[
 F'(z)=\frac1{\pi(1+z^2)}+A
\]
for a constant $A\in\HC$, since the difference has derivative zero on
the connected domain. The integral of $F'$ around any closed curve must
vanish because $F$ is a global primitive. Around a small positive circle
about $\iota$, however, the displayed rational term has integral $1$.
This contradiction independently verifies the failure of same-domain
surjectivity without using the general period theorem.

\subsection{A sphere-average representation of the counterexample}
Let $d\sigma$ be ordinary area measure on the unit imaginary sphere,
with $\sigma(\SSS)=4\pi$, and let
\[
 E(q)=\frac{\bar q}{2\pi^2|q|^4}
\]
be the standard Cauchy--Fueter kernel. The same target admits the
single-valued integral representation
\begin{equation}
 C_b(q)=\frac1{4\pi^2}\int_{\SSS}\frac{d\sigma(J)}{|q-J|^2},
 \qquad
 g_b(q)=2\int_{\SSS}E(q-J)\,d\sigma(J).
 \label{eq:sphereaverage}
\end{equation}
To verify the first formula for $q=x+Ir$, integrate using the angle
between $I$ and $J$:
\[
 \int_{\SSS}\frac{d\sigma(J)}{|q-J|^2}
 =2\pi\int_{-1}^{1}\frac{dt}{x^2+r^2+1-2rt}
 =\frac\pi r\log\frac{x^2+(r+1)^2}{x^2+(r-1)^2}.
\]
This is \eqref{eq:Cb}. If
$\bar\cD=\partial_x-\mathbf i\partial_{x_1}
-\mathbf j\partial_{x_2}-\mathbf k\partial_{x_3}$, then
$g_b=-2\bar\cD C_b$ and
\[
 -2\bar\cD\left(\frac1{4\pi^2|q-J|^2}\right)=2E(q-J).
\]
Differentiation under the integral is valid on every compact set disjoint
from $\SSS$, proving the second formula. It also makes smoothness and
monogenicity away from the sphere transparent: the kernel is monogenic
off its pole. Thus the obstructed target is an averaged Cauchy--Fueter
kernel, not a pathological branchwise construction.

The existence of special logarithmic or inverse-trigonometric Fueter
primitives, including sphere-integrated Cauchy kernels, is not claimed as
new; such kernels occur in the inverse literature~\cite{CPSS}. What the
normalized examples establish here is their unavoidable global period
defect and its two independent coefficients.

\subsection{Which holes count?}
The relevant homology is that of the upper meridian domain $U$, not simply
that of the full symmetric planar domain $D$. For instance, if
\[
 D=\{z:R_1<|z|<R_2\},\qquad 0<R_1<R_2,
\]
then $D$ is not simply connected, but its upper half is parametrized by
$(\log|z|,\arg z)$ with $\arg z\in(0,\pi)$ and is simply connected.
Theorem~\ref{thm:crossreal} therefore gives global inversion for every smooth
axial monogenic function on the four-dimensional shell
$R_1<|q|<R_2$. In contrast, removing a nonreal conjugacy sphere produces
one upper puncture and an eight-real-dimensional cokernel. A blanket rule
using $8\dim H_1(D;\R)$ would give the wrong answer in both comparisons.

\section{Affine residues and the spherical Laurent normal form}
\label{sec:normal}

Fix $p=u+\iota v\in\Cp$, and choose $\varepsilon<v$ so that
$V=B(p,\varepsilon)\subset\Cp$. The circularization of $V\setminus\{p\}$
is a punctured tubular neighborhood of the two-sphere $S_p=u+v\SSS$.
For $q=x+Ir$ in this neighborhood,
\begin{equation}
 \rho=|z-p|=\sqrt{(x-u)^2+(r-v)^2}
             =\operatorname{dist}(q,S_p).
 \label{eq:distance}
\end{equation}
Indeed, among points $u+vJ$ the nearest is obtained with $J=I$.
All growth assertions below are uniform over $I$ and over the planar angle
around $p$.

\begin{definition}[Affine residue pair of a sphere]
For $g\in\cM(V\setminus\{p\})$, define
\begin{equation}
 \resaff_{S_p}g=(a,b)
   :=\left(\int_{|z-p|=\rho}\omega_g,
           \int_{|z-p|=\rho}\theta_g\right),
 \label{eq:resaff}
\end{equation}
where the circle is counterclockwise and $0<\rho<\varepsilon$.
\end{definition}
Closedness and Stokes' theorem on an annulus prove independence of $\rho$.
This is an $\HH^2$-valued invariant of a singular \emph{sphere}. It is not
the usual quaternion-valued three-dimensional flux residue of an isolated
point singularity for the Cauchy--Fueter operator.

\begin{theorem}[Complete spherical singular normal form]
\label{thm:normal}
For every $g\in\cM(V\setminus\{p\})$ there are unique data
\[
 (a,b)\in\HH^2,\qquad c_m\in\HC\ (m\ge1),\qquad g_0\in\cM(V)
\]
such that
\begin{equation}
 g=g_0+G_{p,a,b}
       +\T\left(\sum_{m=1}^{\infty}(z-p)^{-m}c_m\right)
 \quad\text{on }\Omega_{V\setminus\{p\}}.
 \label{eq:normalform}
\end{equation}
The pair is $\resaff_{S_p}g$. The principal-part series and each of its
complex derivatives converge uniformly on compact subannuli. Consequently
its Fueter transform is obtained by termwise application of \eqref{eq:T}.
The function $g$ extends monogenically through the whole sphere if and only
if $a=b=0$ and every $c_m=0$.
\end{theorem}
\begin{proof}
Let $(a,b)$ be the pair in \eqref{eq:resaff}. The function
$g-G_{p,a,b}$ has zero periods on the punctured disk by
\eqref{eq:G-period}. A circle generates its first homology, so
Theorem~\ref{thm:period} supplies a single-valued
$F\in\cS(V\setminus\{p\})$ with
$\T F=g-G_{p,a,b}$.

Apply the scalar Laurent theorem to the four complex components of $F$:
\[
 F=F_0+F_-,\qquad F_0\in\cS(V),\qquad
 F_- =\sum_{m\ge1}(z-p)^{-m}c_m.
\]
The Laurent theorem and the usual holomorphic derivative estimates give
normal convergence of the series and all its derivatives on compact
subannuli. Set $g_0=\T F_0$. On a compact subannulus the factors $1/r$ and
$1/r^2$ are bounded, so \eqref{eq:T} permits termwise transformation and
proves \eqref{eq:normalform}.

The primitive $F$ is unique up to $za+b'$ with quaternionic constants;
these additions have no negative Laurent coefficients and their transforms
vanish. Hence the resulting $c_m$ and $g_0$ are independent of the chosen
primitive. More directly, in two proposed normal forms the residue pairs
agree by integration. The difference of the principal parts then has a
Fueter transform extending to $V$. That extending target has a holomorphic
Fueter primitive $H$ on the disk $V$, by Theorem~\ref{thm:period}. The
principal-part difference minus $H$ has transform zero on the punctured
disk, and is therefore real-affine. Its negative Laurent coefficients all
vanish, proving uniqueness of the principal part and then of $g_0$.

A target extending across the sphere has zero periods on shrinking circles
and admits a local holomorphic primitive on $V$; uniqueness gives zero
singular data. Conversely zero singular data in \eqref{eq:normalform} leave
only the extending term $g_0$.
\end{proof}

\begin{remark}[What the normal form classifies]
This is a classification of \emph{axial monogenic} singularities along a
nonreal two-sphere, expressed in terms of their slice-regular primitives.
It is not a claim to introduce quaternionic Laurent expansions in general.
Slice-regular singularities are treated systematically in~\cite{GPS}, and
Laurent expansions in a related slice Fueter-regular theory occur
in~\cite[\S2.5]{Ghiloni}. The extra datum here is the affine-logarithmic
monodromy that must be removed before a single-valued primitive has an
ordinary complex Laurent expansion.
\end{remark}

\section{Finite growth: a complete classification and sharp removability}
\label{sec:growth}

\begin{theorem}[Finite-growth singular classes]
\label{thm:growth}
Let $N\ge1$ be an integer. For $g\in\cM(V\setminus\{p\})$, the following
are equivalent:
\begin{enumerate}[label=\textup{(\roman*)}]
\item $|g(q)|=O(\rho^{-N})$ uniformly near $S_p$.
\item Its normal form is
\begin{equation}
 g=g_0+G_{p,a,b}
     +\T\left(\sum_{m=1}^{N-1}(z-p)^{-m}c_m\right),
 \label{eq:finiteform}
\end{equation}
with $g_0$ extending through the sphere. The sum is empty for $N=1$.
\end{enumerate}
Modulo monogenic germs extending through $S_p$, the space of such germs is
isomorphic as a real vector space to
\begin{equation}
 \HH^2\oplus(\HC)^{N-1}.
 \label{eq:germspace}
\end{equation}
Its real dimension is $8N$.
\end{theorem}
\begin{proof}
Shrink $V$ if necessary so that $r\ge v/2$. The explicit formulas for the
generator give
\begin{equation}
 |G_{p,a,b}(q)|=O(\rho^{-1}).
 \label{eq:Ggrowth}
\end{equation}
In detail, the numerators $x-u,r-v$ in the $R_p^{-1}$ terms are
$O(\rho)$, while $r$ and $xa+b$ remain bounded and
$s_p=\log\rho$. Since $|\log\rho|=O(\rho^{-1})$, the stated bound follows.

Assume (i), subtract this generator with the actual residue pair, and call
the zero-period remainder $h$. Then $h=O(\rho^{-N})$. Choose its
single-valued primitive $F=zC+E$. Fix an outer circle
$|z-p|=\rho_0$. For an arbitrary inner point, integrate first along the
outer circle and then radially inward. The outer contribution is uniformly
bounded. The one-form estimates give, uniformly in the angle,
\begin{equation}
 |C(z)|+|E(z)|\le
 \begin{cases}
  C_1(1+|\log\rho|),&N=1,\\
  C_N\rho^{1-N},&N\ge2,
 \end{cases}
 \label{eq:primitivegrowth}
\end{equation}
for suitable constants. The same estimates hold for the Euclidean
complexified norm of $F$, because $z$ is bounded on $V$.

For each negative Laurent coefficient, the Cauchy coefficient estimate is
\[
 |c_m|_{\HC}\le \rho^m\max_{|z-p|=\rho}|F(z)|_{\HC}.
\]
If $N=1$, the right side tends to zero for every $m\ge1$; if $N\ge2$,
it tends to zero for every $m\ge N$. Thus only the coefficients allowed in
\eqref{eq:finiteform} can survive.

Conversely, for $F_m=(z-p)^{-m}c_m$, both stem components are
$O(\rho^{-m})$ and their first derivatives are $O(\rho^{-m-1})$.
Equation \eqref{eq:T}, with $r$ bounded away from zero, gives
$\T F_m=O(\rho^{-m-1})$. Together with \eqref{eq:Ggrowth} and the bounded
regular term, this proves (i) from (ii).

The data map to singular germs linearly. The uniqueness and removability
statements of Theorem~\ref{thm:normal} prove injectivity modulo extending
germs, and the displayed finite expression proves surjectivity.
Finally $\dim_\R\HH^2=8$ and $\dim_\R\HC=8$, so the dimension is
$8+8(N-1)=8N$.
\end{proof}

\begin{corollary}[Critical growth and removability]
\label{cor:removable}
For an axial monogenic function near a missing nonreal sphere:
\begin{enumerate}[label=\textup{(\roman*)}]
\item If $g=O(\rho^{-1})$, then $g$ is removable if and only if
$\resaff_{S_p}g=(0,0)$.
\item If $g=o(\rho^{-1})$ uniformly, then $g$ is removable.
\item The little-$o$ condition cannot be weakened to an unrestricted
$O(\rho^{-1})$ condition.
\end{enumerate}
\end{corollary}
\begin{proof}
The first assertion is the case $N=1$ of Theorem~\ref{thm:growth} and the
uniqueness of the normal form. For the second, the circle has length
$2\pi\rho$ and $|z|$ is uniformly bounded there. The estimates
\eqref{eq:periodbounds} make both constant periods tend to zero, so they
vanish. Part (i) applies. For (iii), a nonzero period pair in
$G_{p,a,b}$ gives a nonremovable $O(\rho^{-1})$ example.
\end{proof}

\begin{example}[Zero affine residues do not control higher singularities]
\label{ex:zeroperiodpole}
For $0\ne c\in\HC$, the function
\[
 g=\T\bigl((z-p)^{-1}c\bigr)
\]
has both periods zero and is $O(\rho^{-2})$, but is not removable.
Otherwise a regular local target primitive could be subtracted from
$(z-p)^{-1}c$ to yield an affine function, contradicting its nonzero
Laurent principal part. Thus a residue-only criterion beyond the critical
exponent would be false. Likewise, the transform of
$\exp((z-p)^{-1})$ has zero periods and a nonremovable singularity with
no finite power bound: its primitive has infinitely many negative Laurent
coefficients.
\end{example}

\subsection{Comparison with isolated point singularities}
The exponent $1$ above reflects a singular set of real codimension two and
the axial structure. It is not the exponent for an isolated point in four
real dimensions. For example, the classical Cauchy--Fueter kernel
$\bar q/(2\pi^2|q|^4)$ has size proportional to $|q|^{-3}$; the supplied
article proves an $o(|q|^{-3})$ isolated-point removability criterion
\cite{Source}. Our distance $\rho$ is instead the distance to an entire
nonreal conjugacy sphere. Mixing these two geometries would produce an
incorrect comparison of thresholds or residues.

\section{Exact leading-order laws}
\label{sec:asymptotics}

The preceding classification gives more than bounds: the leading coefficient
has a simple explicit quaternionic norm. For $\alpha,\beta\in\HH$, define
\begin{equation}
 \qnorm(\alpha,\beta)
  =\max_{I\in\SSS}|\beta+I\alpha|
  =\left(|\alpha|^2+|\beta|^2
                +2|\Ima(\beta\bar\alpha)|\right)^{1/2}.
 \label{eq:pairnorm}
\end{equation}
The same maximum is obtained with a minus sign before $I\alpha$.
Indeed $|\beta+I\alpha|^2$ is an affine function of the unit imaginary
vector $I$, with oscillation determined by $\Ima(\beta\bar\alpha)$.
It follows in particular that
\begin{equation}
 (|\alpha|^2+|\beta|^2)^{1/2}
 \le\qnorm(\alpha,\beta)
 \le\sqrt2\,(|\alpha|^2+|\beta|^2)^{1/2}.
 \label{eq:normequi}
\end{equation}
Thus $\qnorm$ is a genuine norm on the eight-real-dimensional pair space.

For a function near $S_p$, write
\[
 M_g(\rho)=\max_{|z-p|=\rho,\ I\in\SSS}|g(x+Ir)|.
\]

\begin{theorem}[Sharp asymptotic constants]
\label{thm:asymptotic}
Let $p=u+\iota v$, $v>0$.
\begin{enumerate}[label=\textup{(\roman*)}]
\item If $g=O(\rho^{-1})$ and its affine residue pair is $(a,b)$, then
\begin{equation}
 \lim_{\rho\downarrow0}\rho M_g(\rho)
   =\frac1\pi\qnorm\left(a,\frac{ua+b}{v}\right).
 \label{eq:criticalconstant}
\end{equation}
\item Suppose the finite normal form of $g$ has highest nonzero Laurent
coefficient $c_m=\alpha+\iota\beta$, with $m\ge1$. Then
\begin{equation}
 \lim_{\rho\downarrow0}\rho^{m+1}M_g(\rho)
   =\frac{2m}{v}\qnorm(\alpha,\beta)>0.
 \label{eq:poleconstant}
\end{equation}
In particular, the transform loses exactly one power of distance at a
nonreal sphere; a nonzero highest coefficient cannot cancel to a smaller
uniform singular order.
\end{enumerate}
\end{theorem}
\begin{proof}
Write $z-p=\rho(\cos t+\iota\sin t)$ and put
$\beta_0=(ua+b)/v$. In \eqref{eq:G-P}--\eqref{eq:G-Q}, expand
$r=v+\rho\sin t$ and $xa+b=ua+b+\rho\cos t\,a$.
Uniformly in $t$ and $I$,
\begin{align}
 P_{p,a,b}
   &=\frac{\sin t\,a+\cos t\,\beta_0}{\pi\rho}
                       +O(1+|\log\rho|),\label{eq:Pasymp}\\
 Q_{p,a,b}
   &=\frac{\cos t\,a-\sin t\,\beta_0}{\pi\rho}
                       +O(1+|\log\rho|).
 \label{eq:Qasymp}
\end{align}
Equivalently,
\begin{equation}
 G_{p,a,b}(x+Ir)
   =\frac{e^{-It}}{\pi\rho}(\beta_0+Ia)
                         +O(1+|\log\rho|).
 \label{eq:Gasymp}
\end{equation}
The unit quaternion $e^{-It}$ preserves norm. In case (i) the rest of the
normal form is regular, by Theorem~\ref{thm:growth}. Multiply
\eqref{eq:Gasymp} by $\rho$, take the supremum, and use uniform convergence
of the error to zero. This proves \eqref{eq:criticalconstant}.

For (ii), let $w=z-p$ and $F_m=w^{-m}(\alpha+\iota\beta)$.
With $\phi=(m+1)t$, the real and imaginary components of
$F_m'=-m w^{-m-1}(\alpha+\iota\beta)$ are
\[
 A_x=-m\rho^{-m-1}(\cos\phi\,\alpha+\sin\phi\,\beta),\qquad
 B_x=-m\rho^{-m-1}(\cos\phi\,\beta-\sin\phi\,\alpha).
\]
Use $A_r=-B_x$, $B_r=A_x$, and $r=v+O(\rho)$ in
\eqref{eq:T}. The result is
\begin{equation}
 \T F_m(x+Ir)=\frac{2m}{v\rho^{m+1}}
          e^{-I(m+1)t}(\beta-I\alpha)+O(\rho^{-m}),
 \label{eq:highestterm}
\end{equation}
uniformly in $t,I$. Lower Laurent terms contribute $O(\rho^{-m})$;
the logarithmic generator is $O(\rho^{-1})$, which is also
$O(\rho^{-m})$ for $m\ge1$; the regular term is bounded.
Taking the rescaled maximum gives \eqref{eq:poleconstant}. Positivity
follows from \eqref{eq:normequi} and $c_m\ne0$.
\end{proof}

The combination $ua+b$ in \eqref{eq:criticalconstant} is not arbitrary.
Proposition~\ref{prop:translation} shows that it is the translation-invariant
constant part of the affine monodromy measured at the real center of the
sphere. The limit is therefore unchanged by a real translation of the
coordinate system, as required geometrically.

\begin{corollary}[An explicit basis and exact orders]
\label{cor:basis}
Let $e_0=1,e_1=\mathbf i,e_2=\mathbf j,e_3=\mathbf k$.
A real basis of the $O(\rho^{-N})$ singular classes is
\[
 G_{p,e_\ell,0},\quad G_{p,0,e_\ell}\quad(0\le\ell\le3),
\]
together with
\[
 \T((z-p)^{-m}e_\ell),\quad
 \T((z-p)^{-m}\iota e_\ell)
 \quad(1\le m\le N-1,\ 0\le\ell\le3).
\]
Every nonzero critical residue class has exact uniform order $\rho^{-1}$;
every class whose highest Laurent index is $m$ has exact uniform order
$\rho^{-m-1}$.
\end{corollary}
\begin{proof}
The basis follows from the data isomorphism \eqref{eq:germspace}.
The exact orders follow from the positive limits in
Theorem~\ref{thm:asymptotic}.
\end{proof}

\section{A holomorphic second-derivative encoding}
\label{sec:encoding}

The period theory has a second description useful for prescribing singularities.
Even if a Fueter primitive is multivalued, its affine monodromy disappears
after two complex derivatives.

\begin{proposition}[The global holomorphic invariant]
\label{prop:encoding}
For $g\in\cM(U)$, choose local holomorphic stems $F$ with $\T F=g$.
Their second derivatives agree on overlaps and define
\[
 Jg\in\mathcal O(U,\HC),\qquad Jg=F''.
\]
For every loop $\gamma$,
\begin{equation}
 a_\gamma(g)=\int_\gamma Jg(z)\dd z,\qquad
 b_\gamma(g)=-\int_\gamma zJg(z)\dd z.
 \label{eq:momentperiods}
\end{equation}
Conversely, a holomorphic $H:U\to\HC$ equals $Jg$ for some
$g\in\cM(U)$ if and only if
\begin{equation}
 \int_\gamma H(z)\dd z\in\HH,\qquad
 -\int_\gamma zH(z)\dd z\in\HH
 \quad\text{for every loop }\gamma.
 \label{eq:realityperiods}
\end{equation}
The kernel of $J$ is the eight-real-dimensional space
\begin{equation}
 \left\{-\frac{2a}{r}-\frac{2I(xa+b)}{r^2}:a,b\in\HH\right\}.
 \label{eq:kernelJ}
\end{equation}
\end{proposition}
\begin{proof}
Local primitives exist on disks by Theorem~\ref{thm:period}. Their
differences are real-affine on every connected overlap, by
Corollary~\ref{cor:kernel}; their second derivatives therefore agree.
On the universal cover, continue a primitive and put
\[
 U_1=F',\qquad U_2=F-zF'.
\]
Then $dU_1=Jg(z)dz$ and $dU_2=-zJg(z)dz$. If the continuation of $F$
changes by $za_\gamma+b_\gamma$, these two functions change respectively
by $a_\gamma$ and $b_\gamma$. This proves \eqref{eq:momentperiods} and
the necessity of \eqref{eq:realityperiods}.

Conversely, integrate $H(z)dz$ and $-zH(z)dz$ on the universal cover to
obtain $U_1,U_2$, and put $F=zU_1+U_2$. Direct differentiation gives
$F'=U_1$ and $F''=H$. Its monodromy is $za_\gamma+b_\gamma$, where the
two coefficients are the integrals in \eqref{eq:realityperiods}. These
coefficients belong to $\HH$, so every branch difference is in
$\ker\T$. The local transforms of $F$ descend to a globally single-valued
$g\in\cM(U)$ with $Jg=H$.

If $Jg=0$, a local primitive is complex-affine. Its real-affine part has
zero transform; its remaining part has the form $\iota(za+b)$.
Equation \eqref{eq:imag-affine} gives \eqref{eq:kernelJ}. The coefficients
are unique, since its first component determines $a$ and its second then
determines $b$. These expressions agree across overlaps, so connectedness
makes $a,b$ global constants. The converse is immediate.
\end{proof}

The reality requirement in \eqref{eq:realityperiods} should not be omitted.
An arbitrary complex-affine monodromy does \emph{not} disappear under the
Fueter transform. Also, $J$ is not injective on product domains; its explicit
kernel is a useful safeguard against incorrectly identifying an axial
monogenic function with the second derivative of a primitive alone.

\begin{corollary}[Ordinary complex residues encode the affine pair]
\label{cor:momentresidues}
If $Jg$ is meromorphic near $p$, then
\begin{equation}
 a=2\pi\iota\,\Res_{z=p}(Jg),\qquad
 b=-2\pi\iota\,\Res_{z=p}(zJg).
 \label{eq:momentresidues}
\end{equation}
For a finite spherical normal form with principal part
$R(z)=\sum_{m=1}^{M}(z-p)^{-m}c_m$, the polar part of $Jg$ is
\begin{equation}
 R''(z)+\frac{a}{2\pi\iota(z-p)}
         -\frac{pa+b}{2\pi\iota(z-p)^2}.
 \label{eq:polarH}
\end{equation}
\end{corollary}
\begin{proof}
The first identities follow from \eqref{eq:momentperiods} and the
componentwise complex residue theorem. For the second, twice differentiate
\eqref{eq:logprimitive}:
\begin{equation}
 \left(\frac{za+b}{2\pi\iota}\log(z-p)\right)''
   =\frac{a}{2\pi\iota(z-p)}
      -\frac{pa+b}{2\pi\iota(z-p)^2}.
 \label{eq:logsecond}
\end{equation}
The regular remainder of the normal form has a holomorphic local primitive,
so its image under $J$ has no polar part.
\end{proof}

\section{A Mittag--Leffler theorem with prescribed affine residues}
\label{sec:ML}

We use the classical scalar Mittag--Leffler theorem in the following standard
form: on a planar domain, arbitrary finite Laurent principal parts prescribed
at a discrete, locally finite set are the principal parts of a meromorphic
function with no other poles. Applied to the four complex coordinates, the
same assertion holds for $\HC$-valued functions. No absolute convergence
of the uncorrected sum of principal parts is required.

\begin{theorem}[Realization of locally finite spherical data]
\label{thm:ML}
Let $V\subset\Cp$ be a simply connected open set, and let
$Z=\{p_j\}_{j\in A}\subset V$ be discrete and locally finite. At each
$p_j$, prescribe independently
\[
 a_j,b_j\in\HH,\qquad
 R_j(z)=\sum_{m=1}^{M_j}(z-p_j)^{-m}c_{jm},
 \quad c_{jm}\in\HC,
\]
where $M_j\ge0$ is finite and an empty sum is allowed.
Then there exists an axial monogenic function $g$ on $\Omega_{V\setminus Z}$
whose local normal form at every $S_{p_j}$ has exactly the affine pair
$(a_j,b_j)$ and the Laurent principal part $R_j$, with no additional
singularities in $\Omega_V$.

Two functions with the same data differ by a function in $\cM(V)$.
The constructed function has a single-valued slice-regular Fueter primitive
on $\Omega_{V\setminus Z}$ if and only if every $a_j=b_j=0$.
There are no compatibility conditions between data at distinct spheres
under these hypotheses.
\end{theorem}
\begin{proof}
For each $j$ prescribe the following meromorphic principal part:
\begin{equation}
 H_j(z)=
 \sum_{m=1}^{M_j}m(m+1)(z-p_j)^{-m-2}c_{jm}
 +\frac{a_j}{2\pi\iota(z-p_j)}
 -\frac{p_j a_j+b_j}{2\pi\iota(z-p_j)^2}.
 \label{eq:MLdata}
\end{equation}
By the componentwise Mittag--Leffler theorem there is a meromorphic
$H:V\to\HC$ with precisely these principal parts and no other poles.
For any closed curve in $V\setminus Z$, the residue theorem gives
\begin{align}
 \int_\gamma H(z)\dd z
   &=\sum_j\wind(\gamma,p_j)a_j,\label{eq:MLperiod1}\\
 -\int_\gamma zH(z)\dd z
   &=\sum_j\wind(\gamma,p_j)b_j.
 \label{eq:MLperiod2}
\end{align}
Only finitely many terms are nonzero for a fixed curve. The use of the
residue theorem here is legitimate because $V$ is simply connected and
$Z$ has no accumulation point in $V$.
To verify the two residues directly, $R_j''$ and $zR_j''$ have no simple
pole. The simple-pole coefficient of $H_j$ is $a_j/(2\pi\iota)$.
In $zH_j$ it is
$(p_j a_j-(p_j a_j+b_j))/(2\pi\iota)=-b_j/(2\pi\iota)$.
Thus both contour integrals belong to $\HH$.

Proposition~\ref{prop:encoding} now constructs $g\in\cM(V\setminus Z)$
from $H$: on the universal cover integrate $Hdz$ and $-zHdz$, set
$F=zU_1+U_2$, and descend $\T F$. This also shows that its affine periods
are exactly \eqref{eq:MLperiod1}--\eqref{eq:MLperiod2}.

It remains to verify the complete local data, not just the periods. Near
$p_j$ consider the multivalued holomorphic expression
\[
 F_j=R_j+\frac{za_j+b_j}{2\pi\iota}\log(z-p_j).
\]
By \eqref{eq:logsecond}, $H-F_j''$ extends holomorphically to a disk about
$p_j$. Choose a holomorphic double primitive $K_j$ of that difference on
the disk. On a branch of the punctured disk,
$(F-F_j-K_j)''=0$, so this difference is complex-affine. Its Fueter
transform is regular on the full disk, since the disk stays away from the
real line. Therefore
\[
 g-\T R_j-G_{p_j,a_j,b_j}
\]
extends through $S_{p_j}$, proving that the local normal form has exactly
the desired data. Away from $Z$ the construction is locally a Fueter
transform of a holomorphic stem, and has no additional singularities.

If two targets have the same data, Theorem~\ref{thm:normal} makes their
difference removable at every $p_j$. Local finiteness permits these local
extensions to be combined into a member of $\cM(V)$. Finally, vanishing
of all the prescribed pairs makes both periods in
\eqref{eq:MLperiod1}--\eqref{eq:MLperiod2} zero on every loop, so a global
primitive exists by Theorem~\ref{thm:period}. Conversely a small loop
about each $p_j$ detects its pair, which must vanish for a global primitive.
\end{proof}

\subsection{Why prescribing the second derivative matters}
A formal sum
\[
 \sum_j\left(R_j+\frac{za_j+b_j}{2\pi\iota}\log(z-p_j)\right)
\]
need not converge and cannot be repaired by silently choosing global
logarithms. The second-derivative construction avoids both problems.
The polar data \eqref{eq:MLdata} are single-valued meromorphic data;
Mittag--Leffler supplies the necessary holomorphic corrections; the two
residue moments guarantee that the remaining primitive monodromy lies in
exactly the real-affine kernel of $\T$. This is why arbitrary residue pairs
can coexist without creating unwanted multivalued target functions.

\begin{example}[Independent defects and poles at several spheres]
\label{ex:several}
Take a simply connected upper domain containing distinct $p_1,p_2,p_3$.
One may prescribe a first-period defect $(\mathbf i,0)$ at $p_1$, a
second-period defect $(0,\mathbf j)$ at $p_2$, and a zero-period principal
part $(z-p_3)^{-2}(1+\iota\mathbf k)$ at $p_3$.
The realization theorem gives a single target with exactly these singular
classes. The first two spheres have exact critical order $\rho^{-1}$;
the third has exact order $\rho^{-3}$ by
Theorem~\ref{thm:asymptotic}. After subtracting
$G_{p_1,\mathbf i,0}+G_{p_2,0,\mathbf j}$, the target has a global
slice-regular primitive on the punctured domain, although its third-sphere
singularity remains. Global integrability and removability are therefore
different questions even for a finite explicit configuration.
\end{example}

\section{What has been established, and what has not}
\label{sec:assessment}

\subsection{A precise contribution ledger}
The local Fueter mapping, the affine kernel, the Vekua equations, and local
inverse constructions are established ingredients. We have not renamed them
as new results. The source manuscript~\cite{Source} provides the immediate
local starting point; the broader inverse literature is
\cite{CSS,CPSS,DQ}. The following conclusions go beyond what is established
in that supplied exposition:

\begin{center}
\small
\renewcommand{\arraystretch}{1.15}
\begin{tabularx}{\textwidth}{@{}p{.28\textwidth}X@{}}
\toprule
Result proved here & Content beyond the local inverse\\
\midrule
Global periods and monodromy & Two target-computable forms give an exact
same-domain solvability test, including both independent obstructions.\\
Realized obstruction space & Explicit normalized generators give an $8m$
real-dimensional cokernel and a stable finite-dimensional correction.\\
Real-axis globalization & The criterion survives on genuine slice domains;
$\HH\setminus\SSS$ supplies concrete failures of unrestricted surjectivity.\\
Spherical normal form & An affine residue pair and a Laurent principal part
classify all local singular data, with a unique regular remainder.\\
Sharp finite growth & The $O(\rho^{-N})$ singular space has dimension $8N$;
exact leading constants and critical removability are computed.\\
Simultaneous realization & Arbitrary locally finite finite principal parts
and affine residue pairs are realized through meromorphic second derivatives.\\
\bottomrule
\end{tabularx}
\end{center}

The proofs supply answers to explicitly formulated problems, rather than
merely suggesting conjectures. On the other hand, it would be inaccurate to
call every auxiliary lemma a new discovery or to claim that the entire
package is already peer-reviewed mathematics. The defensible novelty claim
is narrower: these explicit global and spherical-singularity formulations
were developed in this manuscript and were not found in the primary
literature checked during its preparation.

\subsection{Literature audit and limitations of the priority check}
The search was carried out on 21 September 2026. It included combinations
of ``inverse Fueter'', ``global primitive'', ``multiply connected'',
``period obstruction'', ``affine monodromy'', ``spherical singularity'',
and ``Mittag--Leffler'', together with inspection of the sources listed
below. This is a targeted audit, not an exhaustive search of all Clifford,
Vekua, or generalized analytic function literature.

In particular, the rectangular hypothesis and the explicit spherical-kernel
examples in~\cite{CPSS} were checked in the primary preprint. The contour
formulation of~\cite[Theorem~3.6]{DQ} was inspected. The definitions and the
unqualified surjectivity sentence in~\cite[\S4, p.~2539]{KP} were checked
against the publisher's text. Our counterexample addresses that unrestricted
same-domain assertion; no inference about the validity of unrelated results
in those papers is intended. For singularity context,
\cite{GPS} concerns slice-regular singularities, while
\cite[\S2.5]{Ghiloni} develops Laurent expansions in an octonionic slice
Fueter-regular setting, including real-centered spherical annuli. Neither
context should be conflated with the particular affine-logarithmic
classification proved here.

More recent work continues to broaden the Fueter--Sce phenomenon beyond
the two function classes used here; for example, Ghiloni--Stoppato's
preprint~\cite{GSrecent} studies a multi-step, multi-axial extension.
We cite it as contemporary context, not as evidence that all possible
antecedents to the present period formulas have been excluded.

An important residual uncertainty is historical, not an omitted step in a
stated proof: equivalent obstruction or formal-power descriptions may exist
under different terminology in the generalized analytic-function literature.
Independent expert review should examine both the proof details and that
possibility before a publication-level priority claim is made.

\subsection{Boundaries of the theorems}
Every target in the principal results is \emph{axial}. The conclusions do
not classify arbitrary nonaxial Cauchy--Fueter-regular functions. The local
singularity center satisfies $\Ima p>0$; a sphere collapsing to the real
axis is not covered by the estimates, whose constants explicitly contain
$1/v$. The finite-dimensional global cokernel theorem assumes the stated
finite winding basis. The arbitrary-domain criterion remains valid without
that assumption, but the finite generator sum does not describe arbitrary
infinite-connectivity geometry.

The realization theorem allows finite principal parts at each point and a
locally finite family of points. It imposes no global growth bound at the
boundary or at infinity. Adding weighted norm control, boundary conditions,
or a prescribed behavior as $v\downarrow0$ is a different problem. The
explicit formulas and estimates here provide data for such questions, but
no unstated weighted or boundary-value theorem is being claimed.

\section{Conclusion}
The local affine ambiguity of a Fueter primitive becomes an exact global
obstruction: two quaternion-valued periods per independent upper-meridian
cycle. Because both periods are computable directly from the monogenic
target, they can be tested before attempting inversion and corrected by
explicit single-valued functions. Their geometric significance persists on
domains meeting the real axis.

At a missing nonreal sphere the same obstruction becomes an affine residue
pair. Removing its logarithmic model exposes an ordinary holomorphic
Laurent principal part. This separates three features that a single residue
cannot encode: failure of global integrability, finite-order principal
parts, and removability. The resulting $8N$-dimensional finite-growth
classification, sharp norm limits, and simultaneous realization theorem
form a coherent extension of the local inverse rather than a collection of
isolated analogies with complex analysis.

\appendix
\section{Reproducible checks and artifact guide}
\label{app:checks}

The accompanying file \texttt{verify\_results.py} performs exact symbolic
checks of the two Vekua equations for the branch-free generator, of the
closedness of both one-forms, of the affine-logarithmic second-derivative
identity \eqref{eq:logsecond}, and of its two residue moments. The symbolic
coefficients $a,b$ can be treated as real scalars for these checks: all
coefficients multiplying them in \eqref{eq:G-P}--\eqref{eq:G-Q} are real,
so right $\HH$-linearity extends the component calculation to arbitrary
quaternions. The analytic proofs in the body do not depend on the software.

Numerical checks integrate both forms around circles using a periodic
trapezoidal rule with $16\,384$ equally spaced points. The singular point is
$p=0.3+1.7\iota$, and the detecting circle has radius $0.2$. The general
quaternion pair used is
\[
 a=1-2\mathbf i+\tfrac12\mathbf j+3\mathbf k,\qquad
 b=-0.4+\mathbf i+2\mathbf j-\mathbf k.
\]
The following maximum component errors were obtained in the included run:
\begin{center}
\small
\begin{tabular}{@{}lc@{}}
\toprule
Check & Maximum absolute period error\\
\midrule
Generator with pair $(1,0)$ & $9.548\times10^{-15}$\\
Generator with pair $(0,1)$ & $8.882\times10^{-16}$\\
General quaternion pair & $6.217\times10^{-15}$\\
Reflected generator & $2.620\times10^{-14}$\\
Circle of winding number zero & $2.359\times10^{-15}$\\
\bottomrule
\end{tabular}
\end{center}
For the same quaternion pair the theoretical critical norm limit is
$1.611890081993\ldots$. The sampled values of $\rho M_G(\rho)$ were
\begin{center}
\small
\begin{tabular}{@{}ccc@{}}
\toprule
$\rho$ & Scaled maximum & Absolute difference from the limit\\
\midrule
$10^{-2}$ & $1.635814798836$ & $2.392\times10^{-2}$\\
$10^{-3}$ & $1.615347937083$ & $3.458\times10^{-3}$\\
$10^{-4}$ & $1.612342958506$ & $4.529\times10^{-4}$\\
$10^{-5}$ & $1.611946086722$ & $5.600\times10^{-5}$\\
$10^{-6}$ & $1.611896754419$ & $6.672\times10^{-6}$\\
\bottomrule
\end{tabular}
\end{center}
The maximum over $I\in\SSS$ is evaluated by the exact algebraic formula
\eqref{eq:pairnorm}; only the planar angle is sampled. These numbers are
sanity checks for signs, normalizations, and asymptotic constants, not a
proof by experiment.

The recorded environment was Python~3.13.5, SymPy~1.14.0, and NumPy~2.3.5.
The script prints the versions actually used on each run and raises an
exception if a symbolic identity or a stated numerical tolerance fails.
Run it with
\begin{verbatim}
python verify_results.py
\end{verbatim}
The archive also contains the complete console output in
\texttt{verification.txt}, this editable LaTeX source, its compiled PDF,
and a \texttt{README.md} describing compilation and the status of the work.
There are no external data files or hidden numerical dependencies.

\begin{thebibliography}{10}
\addcontentsline{toc}{section}{References}

\bibitem{Source}
\emph{Classical Complex Theorems over the Quaternions:
Slice-regular and Cauchy--Fueter analogues, with proofs}.
User-supplied unpublished expository manuscript,
\texttt{quaternionic\_analysis.tex}, dated 20 September 2026.
Especially the section ``Fueter's mapping theorem and a constructive local
inverse'', Theorems labelled \texttt{thm:fueter-map} and
\texttt{thm:fueter-inverse}. Used as the starting exposition, not as a
source of priority for the results of this manuscript.

\bibitem{GP}
R.~Ghiloni and A.~Perotti,
\emph{Slice regular functions on real alternative algebras},
Advances in Mathematics \textbf{226} (2011), 1662--1691.
\href{https://doi.org/10.1016/j.aim.2010.08.015}{doi:10.1016/j.aim.2010.08.015}.
\href{https://arxiv.org/abs/1008.4318}{arXiv:1008.4318}.

\bibitem{CSS}
F.~Colombo, I.~Sabadini, and F.~Sommen,
\emph{The inverse Fueter mapping theorem},
Communications on Pure and Applied Analysis \textbf{10} (2011), 1165--1181.
\href{https://doi.org/10.3934/cpaa.2011.10.1165}{doi:10.3934/cpaa.2011.10.1165}.

\bibitem{CPSS}
F.~Colombo, D.~Pe\~na Pe\~na, I.~Sabadini, and F.~Sommen,
\emph{A new integral formula for the inverse Fueter mapping theorem},
Journal of Mathematical Analysis and Applications \textbf{417} (2014), 112--122.
\href{https://arxiv.org/abs/1302.0685}{arXiv:1302.0685}.

\bibitem{DQ}
B.~Dong and T.~Qian,
\emph{Further Aspects of the Fueter Mapping Theorem},
preprint (2018).
\href{https://arxiv.org/abs/1805.02966}{arXiv:1805.02966}.
The primary preprint formulation, especially Theorem~3.6, is the version
consulted here.

\bibitem{KP}
R.~S.~Krau\ss har and A.~Perotti,
\emph{Eigenvalue problems for slice functions},
Annali di Matematica Pura ed Applicata \textbf{201} (2022), 2519--2548.
\href{https://doi.org/10.1007/s10231-022-01208-8}{doi:10.1007/s10231-022-01208-8}.

\bibitem{GPS}
R.~Ghiloni, A.~Perotti, and C.~Stoppato,
\emph{Singularities of slice regular functions over real alternative
$*$-algebras},
Advances in Mathematics \textbf{305} (2017), 1085--1130.
\href{https://doi.org/10.1016/j.aim.2016.10.009}{doi:10.1016/j.aim.2016.10.009}.
\href{https://arxiv.org/abs/1606.03609}{arXiv:1606.03609}.

\bibitem{Ghiloni}
R.~Ghiloni,
\emph{Slice Fueter-Regular Functions},
The Journal of Geometric Analysis \textbf{31} (2021), 11988--12033.
\href{https://doi.org/10.1007/s12220-021-00709-x}{doi:10.1007/s12220-021-00709-x}.
\href{https://arxiv.org/abs/1911.06037}{arXiv:1911.06037}.

\bibitem{GSrecent}
R.~Ghiloni and C.~Stoppato,
\emph{A manifold Fueter--Sce phenomenon in one hypercomplex variable},
preprint (2025).
\href{https://arxiv.org/abs/2511.04771}{arXiv:2511.04771}.

\end{thebibliography}
\end{document}
